\documentclass[12pt]{article}
\usepackage{amsmath, amssymb}
\usepackage[margin=1in]{geometry}
\setlength{\parskip}{6pt}
\title{QUBO Formulation: Mutual-Information Sensor Placement Problem}
\date{}
\begin{document}
\maketitle

\subsection*{Mutual-Information Sensor Placement Problem}

\paragraph{Problem Statement.}
Given a set of $N$ candidate sensor locations, the goal is to select exactly $K$ locations to maximize the total information score while accounting for pairwise redundancy. Each location $i \in \{0, \dots, N-1\}$ is associated with an information score $I_i$, and each distinct pair of locations $\{i, j\}$ has a redundancy penalty $R_{ij}$ that quantifies the loss in mutual information when both sensors are simultaneously active.

\paragraph{Decision Variables.}
For each candidate location $i \in \{0, \dots, N-1\}$, introduce a binary decision variable $x_i \in \{0, 1\}$. The variable $x_i = 1$ indicates that location $i$ is selected, and $x_i = 0$ otherwise.

\paragraph{QUBO Derivation.}
Step 1 --- The original maximization objective is $\max \sum_{i=0}^{N-1} I_i x_i - \sum_{0 \le i < j \le N-1} R_{ij} x_i x_j$. Converting to a minimization problem for the QUBO framework yields:
\begin{align}
  \min \left( -\sum_{i=0}^{N-1} I_i x_i + \sum_{0 \le i < j \le N-1} R_{ij} x_i x_j \right).
\end{align}

Step 2 --- The cardinality constraint requires exactly $K$ sensors to be selected:
\begin{align}
  \sum_{i=0}^{N-1} x_i = K.
\end{align}

Step 3 --- We enforce this constraint via a quadratic penalty term with weight $\lambda > 0$:
\begin{align}
  P(x) &= \lambda \left( \sum_{i=0}^{N-1} x_i - K \right)^2 \\
       &= \lambda \left( \sum_{i=0}^{N-1} x_i^2 - 2K \sum_{i=0}^{N-1} x_i + K^2 + \sum_{i \neq j} x_i x_j \right).
\end{align}
Using the idempotent property $x_i^2 = x_i$ for binary variables, and noting that $\sum_{i \neq j} x_i x_j = 2 \sum_{0 \le i < j \le N-1} x_i x_j$, we expand and collect terms:
\begin{align}
  P(x) &= \lambda \sum_{i=0}^{N-1} (1 - 2K) x_i + 2\lambda \sum_{0 \le i < j \le N-1} x_i x_j + \lambda K^2.
\end{align}

Step 4 --- The total QUBO Hamiltonian $H(x)$ is the sum of the negated objective and the penalty term:
\begin{align}
  H(x) &= -\sum_{i=0}^{N-1} I_i x_i + \sum_{0 \le i < j \le N-1} R_{ij} x_i x_j + \lambda \sum_{i=0}^{N-1} (1 - 2K) x_i + 2\lambda \sum_{0 \le i < j \le N-1} x_i x_j + \lambda K^2.
\end{align}

Step 5 --- Grouping linear and quadratic terms yields the standard QUBO form $H(x) = \sum_{i} Q_{i,i} x_i + \sum_{i < j} Q_{i,j} x_i x_j + c$. The Q matrix entries and offset are given in closed form as:
\begin{align}
  Q_{i,i} &= -I_i + \lambda(1 - 2K), & \text{for } i \in \{0, \dots, N-1\}, \\
  Q_{i,j} &= R_{ij} + 2\lambda, & \text{for } 0 \le i < j \le N-1, \\
  c &= \lambda K^2.
\end{align}
For $i > j$, we set $Q_{i,j} = Q_{j,i}$ to maintain symmetry.

\paragraph{Penalty Weight Justification.}
The penalty weight $\lambda$ must be chosen sufficiently large to ensure that any violation of the cardinality constraint results in an energy strictly greater than the optimal feasible solution. Setting $\lambda > \max_i I_i$ ensures any constraint violation costs more than the maximum possible gain from the objective function. Here $n = N$.

\paragraph{Correctness Argument.}
Minimizing $H(x)$ over $\{0,1\}^N$ recovers the optimal sensor placement because (i) all feasible solutions satisfying $\sum_i x_i = K$ incur zero penalty, so their energy equals the negated objective value; (ii) any infeasible solution violates the cardinality constraint by at least one unit, incurring a penalty of at least $\lambda$. Since $\lambda > \max_i I_i$, the energy of any infeasible solution strictly exceeds that of any feasible solution. Consequently, the global minimum of $H(x)$ must be feasible and corresponds to the configuration that maximizes the original objective.

\paragraph{Illustrative Example.}
Consider a small instance with $N=3$ candidate locations and $K=2$ required sensors. Let the information scores be $I_0 = 10$, $I_1 = 8$, $I_2 = 6$, and the redundancy penalties be $R_{01} = 2$, $R_{02} = 3$, $R_{12} = 1$. We choose $\lambda = 11$, which satisfies $\lambda > \max\{10, 8, 6\}$. Substituting these values into the general formulas yields the concrete Hamiltonian:
\begin{align}
  H(x) &= (-10 - 33)x_0 + (-8 - 33)x_1 + (-6 - 33)x_2 \nonumber \\
       &\quad + (2 + 22)x_0 x_1 + (3 + 22)x_0 x_2 + (1 + 22)x_1 x_2 + 44 \nonumber \\
       &= -43x_0 - 41x_1 - 39x_2 + 24x_0 x_1 + 25x_0 x_2 + 23x_1 x_2 + 44.
\end{align}
Evaluating $H(x)$ over the three feasible assignments gives $H(1,1,0) = -16$, $H(1,0,1) = -13$, and $H(0,1,1) = -13$. The global minimum occurs at $x^* = (1, 1, 0)$ with energy $-16$, corresponding to selecting locations $0$ and $1$ for a maximum total information score of $10 + 8 - 2 = 16$.

\end{document}